\chapter{Pertussis (Whooping Cough)}

\section*{What Is This Test?}
Pertussis testing detects \textit{Bordetella pertussis}, the fastidious gram-negative coccobacillus that causes whooping cough (pertussis). A closely related species, \textit{Bordetella parapertussis}, causes a milder parapertussis syndrome (approximately 5--20\% of pertussis-like illness). Pertussis is a highly contagious respiratory illness characterized by three classic stages: catarrhal (1--2 weeks, indistinguishable from viral URI with rhinorrhea and mild cough), paroxysmal (2--6 weeks, paroxysms of 5--15 forceful coughs during a single exhalation followed by an inspiratory "whoop," post-tussive emesis, and exhaustion), and convalescent (weeks to months, gradual resolution). Infants under 3 months may present with apnea and cyanosis without the classic whoop. Adults and adolescents with waning immunity often present with an undifferentiated prolonged cough (>14 days) without whoop, serving as the primary reservoir for infant transmission. Pertussis is preventable through vaccination (DTaP/Tdap), but immunity wanes 5--10 years after the last dose.

\section*{Alternative Names}
Whooping Cough Test; Bordetella pertussis PCR; Pertussis Culture; Pertussis Serology; B. pertussis NAAT; Bordetella PCR; Pertussis DFA; Tdap Immunity Testing

\section*{Principle}
PCR (NAAT): The preferred diagnostic method. Real-time PCR targeting insertion sequence IS481 (present in multiple copies in \textit{B. pertussis} and in single copies in \textit{B. holmesii} and some \textit{B. parapertussis}) and additional targets (\textit{ptxA}-promoter region, IS1001 for \textit{B. parapertussis}) for specificity. Sensitivity 70--99\%, highest in the catarrhal and early paroxysmal stages (first 2--3 weeks of cough). Specificity >95\%. Results available in 2--4 hours. Culture: Nasopharyngeal specimen inoculated on Regan-Lowe (charcoal-horse blood agar with cephalexin) or Bordet-Gengou agar, incubated at 35--37\textdegree{}C in humidified air for up to 7 days. Colonies are small, shiny, mercury-drop-like. Identification by DFA, biochemical testing, or MALDI-TOF. Culture is 100\% specific but has low sensitivity (20--80\%, lowest in previously vaccinated, antibiotic-treated, or late-stage patients). Direct fluorescent antibody (DFA): Nasopharyngeal smear stained with fluorescein-labeled anti-\textit{B. pertussis} antibody. Sensitivity 50--70\%, specificity variable. Largely replaced by PCR. Serology: ELISA measuring IgG antibodies against pertussis toxin (PT). Sensitivity 80--95\% when tested 2--8 weeks after cough onset. Used for late diagnosis (>2--3 weeks of cough) when PCR and culture sensitivity have declined.

\section*{History}
\textit{Bordetella pertussis} was first isolated by Jules Bordet and Octave Gengou in 1906. The first whole-cell pertussis vaccine (combined with diphtheria and tetanus toxoids as DTP) was introduced in the 1940s, reducing pertussis incidence by >99\% in industrialized countries. Acellular pertussis vaccines (DTaP/Tdap), composed of purified pertussis antigens (pertussis toxin, filamentous hemagglutinin, pertactin, fimbriae types 2/3), replaced whole-cell vaccines in the US starting in 1997 due to fewer local and systemic reactions. However, acellular vaccines induce shorter-lasting immunity and are less effective against \textit{B. pertussis} colonization and transmission compared to whole-cell vaccines. The resurgence of pertussis in highly vaccinated populations beginning in the 2000s (with peaks of >48,000 cases in 2012 in the US) has been attributed to waning acellular vaccine immunity, improved PCR-based diagnosis, and possibly antigenic divergence of circulating strains from vaccine antigens. PCR for pertussis became commercially available in the 2000s and is now the primary diagnostic tool. The CDC recommends Tdap vaccination during each pregnancy (27--36 weeks) to protect newborns through transplacental antibody transfer.

\section*{Why This Test Is Ordered}
\begin{itemize}
  \item Diagnosis of pertussis in any patient with a cough illness lasting 7 days or more with one or more of: paroxysmal coughing, inspiratory whoop, post-tussive emesis, or apnea (in infants <1 year). Testing is most useful in the first 3 weeks of illness when bacterial DNA load is highest
  \item Evaluation of infants under 3 months with apnea, cyanosis, or cough without whoop who are at highest risk for severe and fatal pertussis (particularly before the first DTaP dose at 2 months)
  \item Diagnosis of pertussis in adolescents and adults with prolonged cough (>14 days) where the clinical presentation is atypical and serves as the primary reservoir for transmission to unvaccinated or incompletely vaccinated infants
  \item Investigation of suspected pertussis outbreaks in communities, schools, or healthcare settings: confirmation of the index case and subsequent cases allows appropriate chemoprophylaxis of close contacts and post-exposure vaccination
  \item Confirmation of pertussis as a cause of cough illness in healthcare workers to prevent nosocomial transmission to vulnerable patients (neonates, immunocompromised)
  \item Antibiotic stewardship: differentiate pertussis from other causes of prolonged cough (viral, post-infectious, \textit{Mycoplasma pneumoniae}) to avoid unnecessary antibiotics while ensuring treatment and prophylaxis of confirmed cases
\end{itemize}

\section*{Is This a Primary or Secondary Test}
PCR on nasopharyngeal specimen is the primary diagnostic test for pertussis when the cough duration is less than 3--4 weeks. Culture is a complementary primary test. Serology is a secondary test for patients presenting late (>2--3 weeks of cough) or during outbreaks. DFA is no longer recommended by the CDC due to limited sensitivity and specificity.

\section*{How This Test Is Performed}
Nasopharyngeal (NP) swab or NP aspirate/wash is the required specimen; throat swabs and anterior nasal swabs have unacceptably low sensitivity and should NOT be used. A flexible Dacron or nylon-flocked swab (NOT cotton-tipped, as cotton inhibits PCR and culture) is inserted through one nostril to the posterior nasopharynx, rotated for 5--10 seconds, then withdrawn and placed in Regan-Lowe transport medium (for culture) or manufacturer-specific transport medium (for PCR). For optimal sensitivity, specimens should be collected within 3 weeks of cough onset for PCR and within 2 weeks for culture. Specimens should be transported to the laboratory at room temperature within 24 hours (culture) or per assay requirements (PCR). PCR testing should not be performed on the same specimen used for culture transport medium unless validated, as growth media can interfere with PCR.

\section*{What Preparation Is Needed}
No special patient preparation. Patients should avoid eating or drinking for 30 minutes before nasopharyngeal specimen collection. Antibiotic therapy (macrolides, trimethoprim-sulfamethoxazole) should be withheld until after specimen collection when possible, as antibiotics rapidly reduce organism shedding and culture sensitivity.

\section*{Contraindications and Precautions}
PCR false positives have been reported due to contamination from pertussis vaccine (the vaccine contains pertussis DNA) in the clinic environment or from healthcare worker hands after handling vaccine vials. Clinics should separate vaccine preparation/administration areas from specimen collection areas. False-positive PCR results can also result from cross-reactivity of IS481-targeted PCR with \textit{B. holmesii} (a related organism causing bacteremia in asplenic patients and occasionally pertussis-like illness, but not epidemic pertussis). Laboratories should validate PCR assays that include multiple targets (\textit{ptxA}, IS1001) to distinguish \textit{B. pertussis} from \textit{B. holmesii} and \textit{B. parapertussis}. False-negative PCR results are common after 4 weeks of cough. Culture sensitivity is dramatically reduced by prior antibiotic therapy, vaccination, prolonged transport time, and non-selective media.

\section*{Risks}
Nasopharyngeal swab insertion may cause brief discomfort, sneezing, coughing, or mild epistaxis. In infants with pertussis, the act of nasopharyngeal swabbing may provoke a severe coughing paroxysm with cyanosis or apnea; monitoring during and after the procedure is essential. There is no risk from the laboratory test itself.

\section*{What Happens After the Test}
PCR results available in 2--4 hours (conventional); 1 hour (rapid platforms). Culture results: colonies may appear at 3--7 days; negative cultures held for 7 days. Patients with confirmed pertussis should initiate macrolide antibiotic therapy (azithromycin 500 mg day 1 then 250 mg days 2--5 for adults; or clarithromycin, or trimethoprim-sulfamethoxazole as an alternative) ideally within 3 weeks of cough onset to reduce symptoms; antibiotics after 3 weeks primarily reduce transmission rather than symptoms. The patient should be isolated with droplet precautions until 5 days of antibiotic therapy have been completed (or 21 days after cough onset if not treated). Close contacts should receive post-exposure chemoprophylaxis with the same antibiotic regimen, regardless of vaccination status. The case must be reported to the local public health department (nationally notifiable disease).

\section*{Understanding Results}

\subsection*{Below Normal}
\begin{itemize}
  \item \textbf{PCR negative (not detected):} No \textit{B. pertussis} or \textit{B. parapertussis} DNA detected. May indicate: absence of pertussis infection; viral or other non-pertussis cause of cough; specimen collected >3--4 weeks after cough onset when bacterial DNA levels have declined below the limit of detection; or specimen inadequately collected (anterior nasal or throat swab). Sensitivity of PCR decreases to <50\% after 4 weeks of cough. In an outbreak setting, a negative PCR in an epidemiologically linked symptomatic contact may still represent pertussis and chemoprophylaxis should proceed
  \item \textbf{Culture negative:} No \textit{B. pertussis} isolated. Common in previously vaccinated individuals, patients pretreated with antibiotics, and specimens collected >2 weeks after cough onset, or specimens with delayed transport. Negative culture does NOT exclude pertussis
  \item \textbf{Serology negative (anti-PT IgG < specified cutoff):} Does not support recent pertussis infection. However, serology may be negative in very early infection (first 1--2 weeks of cough) before an IgG response is mounted
\end{itemize}

\subsection*{Above Normal}
\begin{itemize}
  \item \textbf{PCR positive for \textit{B. pertussis} (IS481 and ptxA detected):} Diagnostic of pertussis with high positive predictive value, particularly in a patient with compatible clinical illness. In the absence of clinical symptoms, a positive PCR may reflect asymptomatic carriage or contamination. Results should be interpreted in the context of clinical presentation and epidemiology
  \item \textbf{PCR positive for \textit{B. parapertussis} (IS1001 detected):} Indicates parapertussis, which causes a generally milder but clinically similar illness that may occur in vaccinated individuals; parapertussis is not preventable by pertussis vaccination
  \item \textbf{Culture positive for \textit{B. pertussis}:} Definitive diagnosis (100\% specificity). Allows antimicrobial susceptibility testing (macrolide resistance, though rare in the US, is increasing in China) and molecular typing (PFGE, MLVA, whole-genome sequencing) for outbreak investigations and surveillance
  \item \textbf{Serology positive:} Anti-pertussis toxin IgG level above the cutoff. Single-point serology cutoffs vary by laboratory but are typically $\geq$ 100 IU/mL or $\geq$ 2 standard deviations above the mean of uninfected controls. In older children and adults who have been previously vaccinated, the interpretation is complex; a high anti-PT IgG in the setting of compatible clinical illness (>2 weeks cough) and epidemiologic risk supports the diagnosis. Serology is most useful for late presentations when PCR and culture are negative
\end{itemize}

\section*{Normal Results}
PCR: \textbf{Not detected} for \textit{B. pertussis} and \textit{B. parapertussis}. Culture: \textbf{No \textit{B. pertussis} or \textit{B. parapertussis} isolated at 7 days}. Serology: Anti-pertussis toxin IgG level \textbf{below laboratory-defined cutoff} (typically <40 IU/mL or <20 mcg/mL depending on the assay).

\section*{Follow-up Tests}
\begin{itemize}
  \item \textit{\textbf{B. pertussis}} \textbf{culture (if not already performed):} For PCR-positive cases, culture should be attempted to obtain isolates for antimicrobial susceptibility testing and molecular typing; culture is important for public health surveillance of circulating strains
  \item \textbf{Pertussis serology (anti-PT IgG):} When PCR and culture are negative but clinical suspicion persists (>2 weeks cough), serology may provide diagnostic support
  \item \textbf{Multiplex respiratory pathogen panel:} To evaluate for other causes of prolonged cough (\textit{Mycoplasma pneumoniae}, \textit{Chlamydia pneumoniae}, respiratory viruses) when pertussis testing is negative
  \item \textbf{Chest radiograph:} To evaluate for complications (perihilar infiltrates, atelectasis, pneumothorax, pneumomediastinum) in patients with severe paroxysms or respiratory distress
  \item \textbf{Complete blood count:} Marked leukocytosis (WBC 20,000--50,000 or higher) with lymphocytosis is characteristic of pertussis in the paroxysmal stage, particularly in infants; extreme leukocytosis (WBC >100,000) is associated with poor prognosis and pulmonary hypertension
  \item \textbf{Pharyngeal swab for other Bordetella species:} When PCR identifies \textit{B. holmesii} or \textit{B. parapertussis} to guide clinical management and public health response (parapertussis does not require the same contact investigation)
\end{itemize}

\section*{References}
\begin{enumerate}
  \item CDC. Pertussis (Whooping Cough): Diagnosis and Treatment. Centers for Disease Control and Prevention; 2022.
  \item Mattoo S, Cherry JD. Molecular pathogenesis, epidemiology, and clinical manifestations of respiratory infections due to \textit{Bordetella pertussis} and other \textit{Bordetella} subspecies. Clin Microbiol Rev. 2005;18(2):326-382.
  \item Cherry JD, Tan T, Wirsing von König CH, et al. Clinical definitions of pertussis: summary of a Global Pertussis Initiative roundtable meeting, February 2011. Clin Infect Dis. 2012;54(12):1756-1764.
  \item van der Zee A, Schellekens JF, Mooi FR. Laboratory diagnosis of pertussis. Clin Microbiol Rev. 2015;28(4):1005-1026.
  \item Tiwari T, Murphy TV, Moran J; National Immunization Program, CDC. Recommended antimicrobial agents for the treatment and postexposure prophylaxis of pertussis: 2005 CDC guidelines. MMWR Recomm Rep. 2005;54(RR-14):1-16.
\end{enumerate}

\clearpage
